\documentclass[11pt]{article}
\usepackage[margin=1in]{geometry}
\usepackage{amsmath,amssymb}
\usepackage[T1]{fontenc}
\usepackage[utf8]{inputenc}
\usepackage{lmodern}
\usepackage{enumitem}
\usepackage{hyperref}
\setlist[itemize]{topsep=2pt,itemsep=2pt}
\setlist[enumerate]{topsep=2pt,itemsep=2pt}

\title{Instructions for Counterfactual Search in Hypergraph Classification}
\author{}
\date{}

\begin{document}
\maketitle

\section{Setup}
Let $\theta$ denote the continuous mask parameters optimized by gradient-based search, and let
\[
g(\theta) := -\log p_y(\theta)
\]
be the negative log-likelihood (NLL) of the original predicted class $y$.

The counterfactual objective has the form
\[
\mathcal{L}(\theta;\beta) = -L_{\mathrm{pred}}(\theta) + \beta L_{\mathrm{dist}}(\theta),
\]
where $\beta \ge 0$ controls the trade-off between prediction flipping and sparsity.

A counterfactual is considered \emph{valid} if the final thresholded solution flips the prediction.

\section{1. Incremental Hyperparameter $\beta$ Setting}
The goal is to find the \emph{largest} value of $\beta$ that still yields a valid counterfactual within a fixed trial budget.

\subsection*{Procedure}
\begin{enumerate}
    \item Solve the counterfactual optimization problem with $\beta = 0$.
    \item If no valid counterfactual is found, stop: no admissible $\beta > 0$ can be certified with this procedure.
    \item If a valid counterfactual is found, initialize a small positive value $\beta_{\min}$ (for example $10^{-6}$ or $10^{-4}$).
    \item Increase $\beta$ geometrically (for example by multiplying by $2$) until the first failure or until the trial budget is exhausted.
    \item Once a successful value $\beta_{\mathrm{lo}}$ and a failing value $\beta_{\mathrm{hi}}$ are found, refine the search by binary search in log-space.
    \item Return the largest successful value of $\beta$ found within budget.
\end{enumerate}

\subsection*{Recommended Search Strategy}
Use two phases:
\begin{itemize}
    \item \textbf{Phase A: Exponential bracketing.} Starting from $\beta_{\min}$, test
    \[
    \beta,\; 2\beta,\; 4\beta,\; 8\beta,\; \dots
    \]
    until the first failure is observed.
    \item \textbf{Phase B: Log-binary refinement.} If $\beta_{\mathrm{lo}}$ succeeds and $\beta_{\mathrm{hi}}$ fails, test the geometric midpoint
    \[
    \beta_{\mathrm{mid}} = \sqrt{\beta_{\mathrm{lo}}\beta_{\mathrm{hi}}}
    \]
    and update the interval accordingly.
\end{itemize}

\subsection*{Stopping Rule}
Terminate when either:
\begin{itemize}
    \item the trial budget is exhausted, or
    \item the search interval is sufficiently small for the desired precision.
\end{itemize}

\subsection*{Pseudo-code}
\noindent\textbf{Input:} trial budget $B$, initial small positive value $\beta_{\min}$, geometric factor $\gamma > 1$.\\
\textbf{Assumption:} solving the problem with $\beta=0$ returns a valid counterfactual.

\begin{enumerate}
    \item Run \texttt{SolveCF}$(0)$.
    \item If it fails, \textbf{stop}.
    \item Set:
    \[
    \beta_{\mathrm{best}} \leftarrow 0,
    \qquad
    \beta_{\mathrm{lo}} \leftarrow 0,
    \qquad
    \beta \leftarrow \beta_{\min}.
    \]
    \item \textbf{Exponential bracketing:}
    \begin{enumerate}[label=\alph*)]
        \item While budget remains and \texttt{SolveCF}$(\beta)$ succeeds:
        \begin{itemize}
            \item set $\beta_{\mathrm{best}} \leftarrow \beta$,
            \item set $\beta_{\mathrm{lo}} \leftarrow \beta$,
            \item set $\beta \leftarrow \gamma \beta$.
        \end{itemize}
        \item If a failure occurs, set $\beta_{\mathrm{hi}} \leftarrow \beta$.
    \end{enumerate}
    \item \textbf{Log-binary refinement:}
    \begin{enumerate}[label=\alph*)]
        \item While budget remains:
        \[
        \beta_{\mathrm{mid}} \leftarrow \sqrt{\beta_{\mathrm{lo}}\beta_{\mathrm{hi}}}.
        \]
        \item If \texttt{SolveCF}$(\beta_{\mathrm{mid}})$ succeeds, set
        \[
        \beta_{\mathrm{best}} \leftarrow \beta_{\mathrm{mid}},
        \qquad
        \beta_{\mathrm{lo}} \leftarrow \beta_{\mathrm{mid}}.
        \]
        \item Otherwise set
        \[
        \beta_{\mathrm{hi}} \leftarrow \beta_{\mathrm{mid}}.
        \]
    \end{enumerate}
    \item Return $\beta_{\mathrm{best}}$ and its corresponding valid counterfactual.
\end{enumerate}

\subsection*{Practical Notes}
\begin{itemize}
    \item Warm-start each run from the last successful mask to reduce optimization time.
    \item If optimization is noisy, use a few restarts for each tested $\beta$ and mark success if at least one run flips the prediction.
    \item The primary selection criterion is validity; among valid solutions, maximize $\beta$.
\end{itemize}

\section{2. Dynamic Learning Rate Setting}
Instead of fixing the same learning rate for all datapoints, set it once at the beginning of the optimization from the local gradient at initialization.

\subsection*{Objective for $\beta = 0$}
When $\beta = 0$, the optimization aims only to increase the NLL of the original class:
\[
g(\theta) = -\log p_y(\theta).
\]
Assume a gradient ascent update
\[
\theta_{1} = \theta_0 + \eta \nabla g(\theta_0).
\]
Using a first-order Taylor expansion,
\[
g(\theta_1) - g(\theta_0) \approx \eta \|\nabla g(\theta_0)\|_2^2.
\]

\subsection*{Target Per-Epoch Increase}
To target the initial NLL increase
\[
\Delta = \frac{0.1 + \log N}{\texttt{num\_epochs}},
\]
set the learning rate at initialization to
\[
\boxed{
\eta_0 = \frac{\Delta}{\|\nabla_\theta g(\theta_0)\|_2^2 + \varepsilon}
}
\]
with a small numerical constant $\varepsilon > 0$ (for example $10^{-8}$).

Equivalently,
\[
\boxed{
\eta_0 = 
\frac{(0.1 + \log N)/\texttt{num\_epochs}}
{\|\nabla_\theta(-\log p_y(\theta_0))\|_2^2 + \varepsilon}
}
\]

Then keep the learning rate fixed for all epochs:
\[
\eta \leftarrow \eta_0.
\]

\subsection*{Interpretation}
This choice makes the \emph{first} NLL increase approximately equal to
\[
\frac{0.1 + \log N}{\texttt{num\_epochs}}.
\]
After the first step, the same increase is not guaranteed, because the gradient changes along the optimization path. The rule is therefore a \emph{point-dependent initialization} of the learning rate, not a guarantee of constant increase at every epoch.

\section{Summary}
Use the following two-stage strategy:
\begin{enumerate}
    \item For each datapoint, set the learning rate from the initial gradient using
    \[
    \eta_0 = \frac{(0.1 + \log N)/\texttt{num\_epochs}}
    {\|\nabla_\theta(-\log p_y(\theta_0))\|_2^2 + \varepsilon},
    \]
    and solve the problem with $\beta = 0$.
    \item If a valid counterfactual is found, search for the largest admissible $\beta$ with exponential bracketing followed by log-binary refinement.
\end{enumerate}

\end{document}
